\documentclass{article}
\usepackage[margin=2cm]{geometry}
\usepackage{amsmath}
\usepackage{braidtex}
\begin{document}
\section*{Composable BraidTeX notation}

A primitive diagram with a frame, connector, axis ticks, a region label, and a non-topological dotted path:
\[
  \braid[scale=0.9,frame_top=true,frame_bottom=true,frame_left=true,frame_right=true,frame_axis_color=black,frame_axis_width=0.8]{
    \btin{a b c}
    \btcross{1}\btcross{-2}
    \btframe
    \btticks{topports}{1,2,3}
    \btsep[style={densely dotted},line_width=0.45,tick_ports={1,3}]{1}
    \btconnector[color=black,line_width=0.8,node_radius=0.045,nodes=true,layer=back]{1}{3}{2}
    \btpath{0.15:1.7;0.28:1.0;0.18:0.2}{densely dotted}
    \btregion{0}{2}{u}
  }
\]

Vertical stacking is a model-level operation. The separator is distinct from both crossings and horizontal connectors:
\[
  \btstack[separators=true,labels={u|g},check_labels=false]{
    \braid{\btin{a b c}\btcross{1}\btcross{-2}\btframe}
    \btsep
    \braid{\btin{a b c}\btcross{-1}\btcross{2}\btconnector[color=black,line_width=0.8,node_radius=0.045,nodes=true,layer=back]{1}{3}{2}\btframe}
  }
\]

Operadic substitution replaces one selected input strand with an inserted braid. Since the outer connector touches that strand, \texttt{connector=all} requests separate attachment alternatives:
\[
  \btsubbraid[position=1,connector=all]{
    \btouter{\braid{\btin{a b}\btcross{1}\btconnector{1}{2}{0}}}
    \btinner{\braid{\btin{x_1 x_2}\btcross{-1}}}
  }
\]

Formal diagram sums can be written independently and then composed with other expressions:
\[
  \btalternatives[operator=-]{
    \braid{\btin{a b}\btcross{1}\btout{b a}}
    \btplus
    \braid{\btin{a b}\btcross{-1}\btout{b a}}
  }.
\]
\end{document}
